% !TEX program = xelatex
% 从张量到并矢：新旧对比册（内容修订版）
% 输入（只读）：inputs/old-handout.pdf（20261105 文件版本，29 页）
%               round-04/handout.pdf（48 页，sha256 6e78...3053）
% 修订依据：reviews/tensor-dyad-20261111/comparison/revision-index.json
%           活跃项 COMP-C1..C4、CMP-MATH-01；逐项回应见 issue-response.json
% 生成：python -B render_crops.py  ->  figs/*.png + cropmeta.tex
% 编译：xelatex 两遍
\documentclass[11pt]{ctexart}
%=====================================================================
% theme-violet.tex —— 讲义精讲 · 紫花主题
% 用法：\documentclass[11pt]{ctexart}
%       \input{assets/theme-violet}
%       \begin{document} ... \end{document}
% 只用 ctexart 自带能力 + xcolor/enumitem/fancyhdr/pifont/amsmath，
% 不依赖 tikz/tcolorbox，保证 MiKTeX 可直接编译。
%=====================================================================

\usepackage[a4paper,left=2.0cm,right=2.0cm,top=2.1cm,bottom=2.3cm,%
            headsep=12pt,footskip=22pt]{geometry}
\usepackage{amsmath,amssymb}
\usepackage{needspace}
\usepackage{xcolor}
\usepackage{array}
\usepackage{enumitem}
\usepackage{pifont}
\usepackage{fancyhdr}
\usepackage[hidelinks]{hyperref}

%---------------------------------------------------------------- 配色
\definecolor{violet}{RGB}{107,78,143}      % 主色：紫
\definecolor{violetdeep}{RGB}{63,44,92}    % 深紫：标题/强调
\definecolor{petal}{RGB}{183,159,208}      % 花瓣：装饰线、花饰
\definecolor{bandviolet}{RGB}{247,243,251} % 浅紫底：档案/提示框
\definecolor{bandcream}{RGB}{252,249,242}  % 米白底：易错/注意框
\definecolor{rulepurple}{RGB}{217,207,230} % 分隔线
\definecolor{inkgray}{RGB}{92,96,110}      % 次要文字

%---------------------------------------------------------------- 字体
% 有 Noto CJK 用 Noto，没有就退回 ctex 默认（Windows 下为宋体/黑体）
\IfFontExistsTF{Noto Serif CJK SC}{\setCJKmainfont{Noto Serif CJK SC}}{}
\IfFontExistsTF{Noto Sans CJK SC}{\setCJKsansfont{Noto Sans CJK SC}}{}

%---------------------------------------------------------------- 版式
\setlist{nosep,leftmargin=1.9em}
\linespread{1.28}
\setlength{\parskip}{3.5pt}
\setlength{\parindent}{0pt}
\setlength{\tabcolsep}{5pt}
\renewcommand{\arraystretch}{1.28}

%---------------------------------------------------------------- 花饰
% \ding{95} = 五瓣小花，\ding{96} = 星形花；都来自 pifont，必定存在
\newcommand{\orn}{\textcolor{petal}{\ding{95}}}          % 行内小标
\newcommand{\ornstar}{\textcolor{petal}{\ding{96}}}      % 星形花
\newcommand{\flowerrule}{%                                % 居中花饰分隔线
  \par\vspace{7pt}%
  \noindent\makebox[\linewidth]{\color{petal}%
    \rule{0.20\linewidth}{0.6pt}\hspace{8pt}{\small\ding{95}}\hspace{4pt}%
    \raisebox{1.1pt}{\tiny\ding{95}}\hspace{4pt}{\small\ding{95}}\hspace{8pt}%
    \rule{0.20\linewidth}{0.6pt}}%
  \par\vspace{7pt}}

%---------------------------------------------------------------- 标题与节头
\newcommand{\doctitle}[3]{%   #1 主标题  #2 副标题  #3 元信息（可用 \\ 分行）
  \begin{center}
    {\color{petal}\small \orn\enspace #2\enspace \ornstar}\\[7pt]
    {\LARGE\bfseries\color{violetdeep} #1}\\[10pt]
    {\color{petal}\rule{0.55\linewidth}{1pt}}\\[4pt]
    {\small\color{inkgray} #3}
  \end{center}\vspace{2pt}}

\newcommand{\hdr}[2]{%   #1 章节号（如 §0）  #2 标题
  \par\vspace{12pt}%
  \noindent\colorbox{violetdeep}{\color{white}\bfseries\small\;#1\;}\hspace{7pt}%
  {\color{violetdeep}\large\bfseries #2}\par\nobreak
  \vspace{3pt}{\color{petal}\hrule height 1.2pt}\vspace{6pt}}

\newcommand{\subhdr}[1]{%
  \par\vspace{6pt}%
  \noindent{\color{violet}\bfseries #1}\par\vspace{1pt}}

%---------------------------------------------------------------- 提示框
\newcommand{\callout}[1]{% 浅紫底：讲义档案、读法说明
  \par\medskip\noindent
  \fcolorbox{rulepurple}{bandviolet}{\parbox{\dimexpr\linewidth-2\fboxsep-2\fboxrule}{\vspace{2pt}#1\vspace{2pt}}}%
  \par\medskip}

\newcommand{\warnbox}[1]{% 米白底：易错、边界条件
  \par\medskip\noindent
  \fcolorbox{rulepurple}{bandcream}{\parbox{\dimexpr\linewidth-2\fboxsep-2\fboxrule}{\vspace{2pt}#1\vspace{2pt}}}%
  \par\medskip}

\newcommand{\gapitem}[2]{% #1 缺口类型（紫） #2 内容
  \par\vspace{2pt}\noindent
  \hangindent=1.4em\hangafter=1
  {\color{petal}\small\ding{96}}\enspace{\bfseries\color{violetdeep}#1}\enspace #2\par}

% 题干块（\qstem）：每道例题／作业的「详解」开头必须先复述题干与出处。
% 起因：2026-10-25 第五章讲义把七道作业的题面集中放在一页、详解另起一页，
% 读者翻到「第 4 题」时看不到条件，等于一道没有题干的题目。
% 规则：详解块内自带题干，读者不翻页就能读懂题在问什么。
\newcommand{\qstem}[1]{%
  \par\medskip\noindent
  \fcolorbox{violet}{bandviolet}{%
    \parbox{\dimexpr\linewidth-2\fboxsep-2\fboxrule}{%
      \vspace{2pt}%
      {\small\bfseries\color{violet}题干}\hspace{0.8em}#1%
      \vspace{2pt}}}%
  \par\medskip}

% v2: 保护题干与开头，长题可跨页，续页补上下文。
\newcommand{\qbreak}{\needspace{6\baselineskip}}
% 独立练习：ID、来源、题干、考点、答案标签（通常 ans:ID）。
\newcommand{\practice}[5]{%
  \qbreak\subhdr{练习 #1}\qstem{（#2）#3}%
  \textbf{考点}：#4\par
  详解见附录第~\pageref{#5}~页。\par\medskip}
% 附录开头：ID、来源、条件摘要；自动产生 ans:ID 标签。
\newcommand{\answerstart}[3]{%
  \qbreak\subhdr{#1\quad 详解}\label{ans:#1}\qstem{（#2）#3}}
\newcommand{\answercontinue}[2]{%
  \needspace{5\baselineskip}\subhdr{#1\quad 续解}%
  {\small\color{inkgray}本页接续：#2}\par\medskip}
% 延迟提供默认值，让旧讲义可在导言区保留自己的 \newcommand{\step}。
\AtBeginDocument{\providecommand{\step}[1]{\needspace{3\baselineskip}\par\medskip\textbf{#1}\par}}

% 前提旁注：只放在「需要理解的那一步」后面，给陈述 + 本题代入，不重证。
\newcommand{\aside}[1]{%
  \par\vspace{1pt}\noindent
  \hangindent=1.6em\hangafter=1
  {\small\color{inkgray}%
    \hspace*{0.6em}{\color{petal}$\triangleright$}\hspace{2pt}#1}%
  \par\vspace{2pt}}

%---------------------------------------------------------------- 表格
% 本机 MiKTeX 的 colortbl 与 array 版本不匹配（\rowcolor 会报错），
% 故表格统一用「上粗紫线 + 表头 + 中细线」的三线式，不依赖 colortbl。
\newcommand{\tabtop}{\par\vspace{3pt}{\color{violet}\hrule height 1.2pt}\vspace{1pt}}
\newcommand{\tabmid}{\par\vspace{1pt}{\color{petal}\hrule height 0.5pt}\vspace{1pt}}
\newcommand{\tabbot}{\par\vspace{3pt}{\color{rulepurple}\hrule height 0.5pt}\vspace{8pt}}
\newcommand{\thead}[1]{{\bfseries\color{violetdeep}#1}}

%---------------------------------------------------------------- 页眉页脚
\pagestyle{fancy}
\fancyhf{}
\renewcommand{\headrulewidth}{0.4pt}
\renewcommand{\headrule}{\hbox to\headwidth{\color{rulepurple}\leaders\hrule height \headrulewidth\hfill}}
\fancyhead[L]{\small\color{inkgray}\songti \doctag}
\fancyhead[R]{\small\color{inkgray}\songti \coursetag}
\fancyfoot[C]{\small\color{inkgray}\orn\enspace{\color{violet}\thepage}\enspace\orn}

% 每篇讲义在导言区设这两个值
\newcommand{\doctag}{}
\newcommand{\coursetag}{}
\newcommand{\settags}[2]{\renewcommand{\doctag}{#1}\renewcommand{\coursetag}{#2}}

%---------------------------------------------------------------- 数学记号
\newcommand{\vect}[1]{\boldsymbol{#1}}
\newcommand{\dd}{\,\mathrm{d}}
\newcommand{\ddx}{\frac{\mathrm{d}}{\mathrm{d}x}}
\addtolength{\jot}{2pt}



\usepackage{graphicx}
\usepackage{tabularx}
\usepackage{longtable}
\settags{新旧对比册}{理论力学}
\setlength{\headheight}{22pt}
\setlength{\emergencystretch}{3em}
\graphicspath{{figs/}}

\makeatletter
\newcommand{\cropdata}[5]{\@namedef{crop@#1}{#2 p#3\enspace·\enspace #4\enspace·\enspace #5}}
\newcommand{\cropimg}[2][1.0]{%
  \par\medskip\noindent
  \begin{minipage}{\linewidth}%
  {\small\bfseries\color{violet}\csname crop@#2\endcsname}\par\vspace{2pt}%
  \fcolorbox{rulepurple}{white}{\includegraphics[width=\dimexpr#1\linewidth-2\fboxsep-2\fboxrule\relax]{figs/#2.png}}\par
  \end{minipage}\par\medskip}
\makeatother
% generated by render_crops.py -- do not edit by hand
\cropdata{old-p01-audience}{旧稿}{1}{y=176--262}{含“为谁写”整段读者基础声明}
\cropdata{old-p01-route}{旧稿}{1}{y=706--772}{含“读法”整段与最短路径}
\cropdata{old-p11-summation}{旧稿}{11}{y=300--458}{含§4.3标题与求和约定整段，包括原文给出的规则本身}
\cropdata{old-p11-law}{旧稿}{11}{y=676--792}{含§4.4标题与实际变换律}
\cropdata{old-p12-quotient}{旧稿}{12}{y=326--470}{含§4.5商法则与原证明末句}
\cropdata{old-p15-epsilon}{旧稿}{15}{y=160--330}{含对偶张量式与$\varepsilon$的使用}
\cropdata{old-p15-note}{旧稿}{15}{y=377--432}{含旧稿旁注里已有的A23分量核对，不能裁掉}
\cropdata{old-p17-center}{旧稿}{17}{y=138--320}{含§7.1开头、v=$\omega$×r断言与§5.4引用}
\cropdata{old-p20-symtop}{旧稿}{20}{y=640--792}{含旧§8.3对称陀螺、方程、“标准振子方程”与解的形式}
\cropdata{old-p21-symtop}{旧稿}{21}{y=64--172}{紧接解式之后：只给解释与几何图像，没有初值或代回核对}
\cropdata{old-p21-continuum}{旧稿}{21}{y=394--530}{含§9.1连续体与面力原句}
\cropdata{old-p22-fig5}{旧稿}{22}{y=148--310}{含图5与图注的两端面符号}
\cropdata{old-p22-balance}{旧稿}{22}{y=320--500}{含旧§9.4高斯定理式与平衡方程，只到“写得开”一句}
\cropdata{old-p23-eig}{旧稿}{23}{y=376--430}{含$\lambda$=5的中间步与随后给出的n1}
\cropdata{new-p01-assume}{新稿}{1}{y=180--300}{含“本稿只假设你会”清单与必读/选读读法}
\cropdata{new-p02-actions}{新稿}{2}{y=230--310}{含§0起点六个动作与读者处理方式}
\cropdata{new-p02-route}{新稿}{2}{y=146--216}{含最短主线整段}
\cropdata{new-p03-summation}{新稿}{3}{y=104--192}{含求和约定定义与同一指标出现三次的提醒}
\cropdata{new-p16-epsilon}{新稿}{16}{y=684--794}{含$\varepsilon$的六个排列定义与第一分量核对}
\cropdata{new-p17-quotient}{新稿}{17}{y=628--788}{含商法则三步证明与T$'$=RTTR}
\cropdata{new-p23-center}{新稿}{23}{y=410--556}{含v=$\omega$×r解释与质心/瞬时轴旁注}
\cropdata{new-p29-symtop}{新稿}{29}{y=206--392}{含新§5.3标题、方程、二维矢量导数式与“不借助微分方程理论”的入口}
\cropdata{new-p29-bridge}{新稿}{29}{y=395--685}{含点乘推导、直接代回核对与由初始横向矢量定常数}
\cropdata{new-p30-continuum}{新稿}{30}{y=578--634}{含连续体与面力矢量定义}
\cropdata{new-p32-derive}{新稿}{32}{y=286--470}{含新§6.4标题、偏导含义与小立方体净力推导}
\cropdata{new-p32-example}{新稿}{32}{y=478--676}{含应力线性变化的净力例子与均匀应力对照}
\cropdata{new-p32-fig8}{新稿}{32}{y=66--272}{含图8与两端面符号}
\cropdata{new-p38-p1}{新稿}{38}{y=505--700}{含$\lambda$=5行方程、单位化与代回核验}


\newcommand{\oldlbl}{\par\smallskip\noindent{\bfseries\color{violetdeep}旧稿卡在哪里。}\hspace{0.2em}}
\newcommand{\newlbl}{\par\smallskip\noindent{\bfseries\color{violetdeep}新稿改了什么。}\hspace{0.2em}}
\newcommand{\bndlbl}{\par\smallskip\noindent{\bfseries\color{violetdeep}边界与未验证。}\hspace{0.2em}}
\newcommand{\cmpitem}[2]{\needspace{7\baselineskip}\subhdr{#1\quad #2}}
\newcommand{\srcsmall}[1]{{\small\color{inkgray}#1}}

\begin{document}

\doctitle{从张量到并矢：新旧对比册（修订版）}%
{旧稿 20261105 文件版本（29 页）与 round-04 新稿（48 页）的页面级对照}%
{2026-11-11\enspace·\enspace DSH 制作\enspace·\enspace 按内容审核问题单修订，只比较实际成品页面}

\callout{\raggedright\textbf{本次性质与分工。}本次不是从零试生成，也不是协调器代写。以用户指定的旧稿
\texttt{inputs/old-handout.pdf}（冻结的 2026-11-05 文件版本，物理页 1--29；封面日期仍印
2026-11-03，所以按文件版本称 20261105 版，不据封面改称 20261103 版）和 \texttt{inputs} 内原资料为输入，
由 DSH 独立返工出新稿 \texttt{round-04/handout.pdf}（48 页）。协调器只维护 skill、给范围与只读审核，
没有写正文、图、解答或脚本。本册由 DSH 制作并自检。本册是原对比册的内容修订版：只改比较评价、引文与
必要裁剪，不改 \texttt{round-04} 讲义。原对比册的版面已由外部视觉审核接受，但内容审核提出 5 项
（COMP-C1--C4、CMP-MATH-01）；本修订版逐项回应，记录在 \texttt{issue-response.json}，仍需外部复核。}

\begin{center}
\small
\begin{tabularx}{\linewidth}{@{}p{1.35cm}p{4.25cm}p{1.15cm}X@{}}
\thead{角色} & \thead{文件} & \thead{物理页} & \thead{版本与摘要}\\
\hline
旧稿 & \texttt{inputs/old-handout.pdf} & 29 & 文件版本 20261105；封面 20261103；SHA-256 \texttt{97bb...9462}\\
新稿 & \texttt{round-04/handout.pdf} & 48 & skill 2.2.2，round-04；SHA-256 \texttt{6e78...3053}；审核状态 \mbox{accepted}（content/math/visual scope）\\
\end{tabularx}
\end{center}

\warnbox{\textbf{事实边界。}\;(1) 学习者只自述“只会简单的线性代数知识”，未做逐项诊断；本册把这句话当作
\emph{声明}，不当作已核实的掌握清单。(2) 学习者实际试读尚未发生；本册不声称任何人已学会、已读懂，也不声称
教学效果已经验证。(3) 本册不把页数更多、例子更多、自动检查更多当作“效果更好”的证据；它只指出具体页面上的
具体差别。(4) DSH 自身没有图像输入：本册的文字层/坐标检查不能代替看图，本册 28 处裁剪与本册最终页
渲染仍待支持看图的审核者目视。新稿 round-04 的外部内容、数学、视觉审核已完成并记录在
\texttt{final-selection.json} 与 \texttt{visual-review-round04.json}；这里说的是 DSH 没有自己看图，
不是新稿缺少外部视觉审核。文中对“图注文字”的比较只依据 PDF 文字层，不冒充看图。
(5) 旧稿 29 页、新稿 48 页是各自的实际页数，不是评价指标。}

\vspace{4pt}
\srcsmall{引用约定：以下“旧 p11”指旧稿 PDF 的物理第 11 页，“新 p16”指新稿 PDF 的物理第 16 页；
“§4.3”“§2.6”分别是两稿自己的小节号。裁剪矩形以 PDF 点（pt）计，页面 A4 为 595$\times$842 pt。}

\clearpage
\hdr{一}{输入、页码与读法}

本册只回答四个问题：旧稿的简单线性代数起点有没有得到必要支持；具体例图与中间步骤是否真的有用；
最短路线与选读是否真的走得通；新稿补的局部验证/初值与“应力变化$\rightarrow$净力”桥梁是否实际印出，
以及是否引入了没有解释的新术语。每组都选两边都实际存在、能指到具体句子和具体页面的段落，不构造旧稿
的弱点。找不到具体差别的主题不写；原稿中缺乏实际差别的组已合并或替换，不为保留组数制造改善。

\begin{center}
\small
\begin{tabularx}{\linewidth}{@{}p{1.55cm}Xp{2.75cm}@{}}
\thead{组} & \thead{比较内容} & \thead{本册位置}\\
\hline
A & 起点支持：分量、求和约定、$\varepsilon_{ijk}$ 是否在正文可用 & §A\\
B & 中间步骤与例图：商法则证明、例 1 主轴、图 5/图 8 端面符号 & §B\\
C & 选读与主线：最短路线，以及 §5.3 局部验证/初值、§6.4 应力变化与净力桥梁 & §C\\
D & 未解释术语：定点/质心、连续体/面力；瞬时轴是新稿自己引入并解释的术语 & §D\\
\end{tabularx}
\end{center}

\textbf{裁剪原则。}每张裁剪都保留该小节的标题或紧邻的前提句，不裁掉判断所依据的上一句；每张裁剪在原页面上的
矩形和“必须出现在裁剪里的字符串”都写进 \texttt{crops.json}，由 \texttt{render\_crops.py} 逐张核对。渲染 200 dpi，
按页宽（多图项按 0.9 页宽）排入本册，等效约 190--214 dpi；正文 11 pt 原字在纸上约 7.8--8.6 pt，可放大阅读。正文引用原句时直接抄录，
不润色成“更该如此”的样子。

\textbf{每项三段。}“旧稿卡在哪里”只写读者在该页会实际撞到的句子或缺口；“新稿改了什么”只写新稿页面上实际
印出的内容；“边界与未验证”区分已核对的文字/数值与尚未做的目视、试读。

\callout{\textbf{这不是评审结论。}\texttt{final-selection.json} 已接受新稿（accepted within documented
content, math and visual scope；\texttt{visual-review-round04.json} 记录外部视觉审核完成），但同页写明
comparison 当时仍待制作，且 review\_limits 明说“学习改善与用户理解尚未由学习者检验”“接受反映已审内容
与页面，不按页数或断言条数”。本册是新旧页面差异的说明，不替代评审报告；本修订版仍需外部复核。}

\clearpage
\hdr{二}{轮次与规则：真实历史，不写成一次提示}

\begin{center}
\small
\begin{tabularx}{\linewidth}{@{}p{1.15cm}p{1.15cm}p{2.6cm}p{1.05cm}X@{}}
\thead{轮次} & \thead{skill} & \thead{本轮具体问题单} & \thead{页数} & \thead{处理结果}\\
\hline
R1 & 2.2.0 & 无 & 45 & 以旧稿和原资料独立重修；未收到本轮具体问题单\\
R2 & 2.2.1 & 有（16 项） & 47 & DS 按问题单修改\\
R3 & 2.2.1 & 有（5 项） & 48 & DS 按问题单修改\\
R4 & 2.2.2 & 有（1 项 V8） & 48 & DS 按问题单修改；最终接受\\
\end{tabularx}
\end{center}

\textbf{不能含糊的三件事。}第一，R1 是“没有本轮具体问题单”的第一版：它已经以用户指定旧稿为范围与题目来源、
以原资料为依据独立重写，因此新稿的起点支持（§0 六个动作、必读/选读标记、最短主线）在 R1 就出现了，不能记到
后续反馈或 2.2.1/2.2.2 规则头上。第二，R2--R4 都是收到问题单后的 DS 修改，逐项证据在各轮
\texttt{issue-response.json}；这不是一句提示自动完成全部改善。第三，R1 也照搬了旧稿的两处问题——例 1 的中间步
（把 $-e_1+e_2$ 当作 $\lambda=5$ 的方向）和把 $(1,1,1)/\sqrt3$ 叫“体对角线”；前者由 R2 的 M7、后者由 R2 的 M6
修正。可见“第一版自动全对”不成立。

\textbf{规则变化与暴露的问题。}2.2.1 增加三条检查：推荐阅读路线是否真的通到后面的应用/练习；例子是否真的
激活它要说明的现象或术语；输入/槽位/顺序与实际印出的中间推理是否一致，而不只看意图中的最终公式。R2 的 M2（固定槽位与右作用对齐）、M7（例 1 中间步）、S3（例子没有激活它要解释的错误）正好落在这三条上，
但这些同时是具体问题单里的条目，不能把改动单独归因给某条规则。2.2.2 只要求内容/版面改动后复查受影响的下游题块与续解
标记；R4 只移动 P11 的续解位置，没有改内容，所以本册的内容差异不来自 2.2.2。规则修改与外部具体反馈共同产生了
变化，二者不能互相冒充。本册 §C 的 C2、C3 选的正是 round-04 接受范围内已保留的两处阅读桥梁：round-04 的
\texttt{issue-response.json} 记录该轮只改 V8 分页，说明 §5.3 与 §6.4 的讲解在此前轮次已进入并通过内容/阅读复核，
本轮保留。选它们是为了贴近本次学习卡点，不是把“新稿多了内容”当作效果证据。

\srcsmall{依据：\texttt{final-selection.json} 的 generation\_history、rule\_updates 与各轮 issue-response.json。
本册不引用父项目其它讲义，也不读协调器成品脚本。}

\clearpage
\hdr{三}{比较方法：只比较能指到句子的差别}

\begin{itemize}
\item \textbf{选段。}先在两稿文字层里找同一主题，再确认两边都有可引用的句子与页面；只有一边有、或只能靠常识
推断“旧稿大概没讲”的，不写。
\item \textbf{裁剪。}用 PyMuPDF 按矩形裁原页面，保留小节标题与前提句；每张裁剪的矩形、DPI、像素尺寸、SHA-256 和
必须包含的字符串都写进 \texttt{render-crops.json}。裁剪不缩放原字，只按页宽排入。
\item \textbf{数值。}\texttt{verify\_old\_new\_math.py} 复核本册实际引用的数值/代数断言：旧稿商法则末句的
$T'=RTR^{\mathsf T}$ 与它自己 §4.4 的 $T'=R^{\mathsf T}TR$ 不同、且只有后者与逐项投影一致；例 1 的
$(D-5\hat1)(1,1,0)=0$ 而 $(-1,1,0)$ 不是零空间；$\varepsilon$ 符号给出 $A_{23}=-a_1$；图 8 的
$t(\pm e_1)=\pm\sigma_0e_1$ 与 $30^\circ$ 截面的 $t_n=\tfrac34\sigma_0$、$\tau=\tfrac{\sqrt3}{4}\sigma_0$；
对称陀螺 $A=B=4$、$C=5$、$\omega_3=6$ 给 $\Omega=1.5$、$T\approx4.19\,$s，且写出的解代回两条一阶方程；
应力场 $\sigma_{11}=\sigma_0x_1/L$ 给 $(\nabla\cdot\sigma)_1=\sigma_0/L$，均匀应力给零。
\item \textbf{不做的。}不按“新稿多了几张图/几道题/多少条断言”评分；不把审核通过当作学习效果；不对未读的页面
下结论。本册只覆盖实际列出的 28 处裁剪。
\end{itemize}

{\raggedright
\textbf{复现命令。}在 \texttt{comparison-revised/} 下：\texttt{python -B render\_crops.py} 生成裁剪与
\texttt{cropmeta.tex}；\texttt{xelatex -halt-on-error -interaction=nonstopmode comparison.tex} 两遍；
\texttt{python -B verify\_old\_new\_math.py}；\texttt{python -B check\_compare.py --strict}；
\texttt{python -B verification\_report.py} 生成 \texttt{verification.txt}。最终每页渲染见
\texttt{render/}，也可运行 \texttt{build.ps1} 一次完成。\par}

\clearpage
\hdr{A}{起点支持：简单线性代数能否直接接上}

旧稿第 1 页说，它是“为‘会算矩阵乘法，但二次型、迹、特征值都很生疏，张量完全没学过’的读者写”。这个声明比
本轮自述更具体：它要求会矩阵乘法，并把基、分量、点乘、叉乘都当作已会，或由课件 p8 提供。求和约定不同：旧稿 §4.3（p11）已在正文写出规则，只是位置放在前面几节按分量写开之后。新稿第 1 页
把假设收窄成一张可核对的清单：“写矢量的分量、做点乘与叉乘、做矩阵乘矢量、取转置”，再加“导数与积分各一句
话”；然后用 §0 的六个动作各练一个能口算的例子。

\cropimg{old-p01-audience}
\cropimg{new-p01-assume}
\cropimg{new-p02-actions}

\begin{center}
\small
\begin{tabularx}{\linewidth}{@{}p{3.15cm}p{4.95cm}X@{}}
\thead{六个动作} & \thead{旧稿首次直接使用} & \thead{新稿首次给出/复练}\\
\hline
写分量、点乘、叉乘 & 旧 p2 §1.1 基展开；旧 p3 §1.3 点乘、叉乘 & 新 p2 §0.1 定义基、分量、$\delta_{ik}$、点乘、叉乘，并给口算例\\
矩阵乘矢量 & 旧 p3 §1.2 & 新 p4 §0.3\\
换基 & 旧 p3--4 §1.3 & 新 p4--6 §0.4\\
迹、行列式 & 旧 p5 §2.1--2.2 & 新 p6--7 §0.5\\
\end{tabularx}
\end{center}

\cmpitem{A1}{求和约定：旧稿在 §4.3 写出规则，新稿前移到 §0.1 并配口算例}
\cropimg{old-p11-summation}
\cropimg{new-p03-summation}

\oldlbl 旧稿 §4.3（p11）在正文里写出规则本身：“课件 p8 已经给了求和约定：同一项里重复出现两次的指标
自动求和，不写 $\Sigma$。本讲义正文照旧写求和号（免得看错缩并了哪个指标），只在 §6 的字典与 §7 之后的
推导里用约定。”旧稿附录 A（p28）又重复一次。它不是“约定没有在旧稿正文出现”，而是把定义放在 §4.3，
并在前面几节先按分量写开；同页没有再配一个点乘/矩阵作用的实例。

\newlbl 新稿 §0.1（p3）在第一次使用前写同一条规则：“同一项里同一个指标出现两次，就自动对 1,2,3 求和”，
给 $a\cdot b=a_kb_k$、$(Tx)_i=T_{ij}x_j$、$a_i=\delta_{ij}a_j$ 三例，并提醒同一指标出现三次不是约定用法；
随后用 $a=(1,2,2)$、$b=(2,1,0)$ 的口算例把点乘、投影与叉乘各走一遍。差别是位置前移与操作例，不是旧稿
缺定义。

\bndlbl 这条只保证“能读式子”；点乘、叉乘、矩阵乘矢量仍按清单假定已会，新稿没有重教。旧稿的规则本身
正确，附录 A 也重复；本册不再把“约定没在旧稿正文出现”当作差异。约定段尚未经过学习者试读，是否一次
读懂未知。

\needspace{26\baselineskip}
\cmpitem{A2}{$\varepsilon_{ijk}$：旧稿有局部 $A_{23}$ 核对，新稿补六个排列与重复指标规则}
\cropimg{old-p15-epsilon}
\cropimg{old-p15-note}
\cropimg{new-p16-epsilon}

\oldlbl 旧稿 §5.5（p15）直接写 $A_{ij}=-\varepsilon_{ijk}a_k$、$a_k=-\tfrac12\varepsilon_{kij}A_{ij}$，并说
“就是课件 p8 的对偶张量”；正文没有列全偶/奇排列与重复指标为零的通用表。但旧稿同页下方旁注已经做了
分量核对：“用分量核一遍最省事：第三种写法给出 $A_{23}=-\varepsilon_{231}a_1=-a_1$，与课件矩阵的
$(2,3)$ 元一致。”所以旧稿并非只能背下三个矩阵元；它缺的是可迁移的六排列/重复指标规则，不是
这个例子的核对。

\newlbl 新稿 §2.6（p16）定义 $\varepsilon_{ijk}$：偶排列（123, 231, 312）取 $+1$，奇排列（132, 213, 321）取
$-1$，有重复指标取 0，并说“六个排列已在括号里列全，直接查表”；接着核对第一分量
$\varepsilon_{123}a_2b_3+\varepsilon_{132}a_3b_2=a_2b_3-a_3b_2$。p17 还注明课件 p8 与 RIT Whelan notes 10
的符号约定差别。

\bndlbl 新稿把“偶/奇排列”压在一句括号里，仍要求读者愿意按六个排列查表；这是可操作但仍需动手的第一步，
未试读。旧稿的对偶张量公式本身与它 p15 的矩阵一致；本册不再说旧稿读者无法自己核对 $A_{23}$，只保留
“旧稿没有列全六排列/重复指标规则”这个较窄的差异。

\clearpage
\hdr{B}{中间步骤与例图：实际印出的推理是否自洽}

\cmpitem{B1}{商法则证明：旧稿末句的转置位置与自己的 §4.4 相反}
\cropimg[0.9]{old-p12-quotient}
\cropimg[0.9]{old-p11-law}
\cropimg[0.9]{new-p17-quotient}

\oldlbl 旧稿 §4.5（p12）写：“把 $u'=Ru$、$v'=Rv$ 代进去，与 $u=TRv$ 比较，就逼出 $T'=RTR^{\mathsf T}$，
正是 §4.4 的变换律。”但 §4.4（p11）自己印的是 $T'_{ij}=R_{ki}R_{lj}T_{kl}$，即 $T'=R^{\mathsf T}TR$。
两者不是同一个矩阵；按 p12 的字面走，读者会把 $RTR^{\mathsf T}$ 当成换基后的账本，而旧稿 p4 自己刚提醒过
$RTR^{\mathsf T}$ 是“另一台机器”。

\newlbl 新稿 §2.7（p17）把证明写成三步：$u'=T'v'$，代入 $u'=R^{\mathsf T}u$、$v'=R^{\mathsf T}v$ 得
$R^{\mathsf T}u=T'R^{\mathsf T}v$，于是 $u=RT'R^{\mathsf T}v$；与 $u=Tv$ 比较得 $T=RT'R^{\mathsf T}$，两边
左乘 $R^{\mathsf T}$、右乘 $R$，得 $T'=R^{\mathsf T}TR$。三步都印在页面上。

\bndlbl 这是一处证明书写差异，不推翻旧稿“$\sigma$ 一定是张量”的结论；旧稿的结论仍成立，只是中间一步写反。
新稿的写法已由本册数值断言复核（只有 $R^{\mathsf T}TR$ 与逐项投影一致），但学习者是否跟得上三步未试读。

\cmpitem{B2}{例 1 主轴：旧稿中间步给出的方向不是零空间}
\cropimg{old-p23-eig}
\cropimg{new-p38-p1}

\oldlbl 旧稿 p23 写：“$\lambda=5$ 时 $(D-5\hat1)$ 的前两块给出 $-e_1+e_2$ 方向，即 $n_1=(e_1+e_2)/\sqrt2$。”
前半句给出的 $(-1,1,0)$ 不是 $D-5\hat1$ 的零空间；后半句的 $n_1$ 才是正确的。中间步与最终答案不一致，
且没有代回核验；照着“$-e_1+e_2$”走的读者会得到错误特征向量。

\newlbl 新稿附录 A 的 P1（p38）先印出 $D-5\hat1$，由前两行 $-x+y=0$、第三行 $-3z=0$ 得 $(1,1,0)$，单位化
$n_1=(e_1+e_2)/\sqrt2$，再代回 $(D-5\hat1)(1,1,0)=0$ 核验；$\lambda=3$、$\lambda=2$ 也补了行方程。注意：
这不是第一版就有的——R1 的 P1（p36）与旧稿同样写“前两块给出 $-e_1+e_2$”，R2 按问题单 M7 才改。

\bndlbl 新稿把这道题从正文例题移到附录练习详解，正文只给题干与考点；这是结构选择，不构成“更容易学会”的
证据。旧稿例 1 其余步骤（$D\cdot x$、$\mathrm{tr}$、$\det$、$D:D$、特征方程、谱分解）的数值正确。

\needspace{26\baselineskip}
\cmpitem{B3}{图 5 与图 8：端面符号的文字层差别}
\cropimg{old-p22-fig5}
\cropimg{new-p32-fig8}

\oldlbl 旧稿 p22 图注写“两端面（法向沿 $\pm e_1$）上的面力是 $\sigma_0e_1$”。按 $t(n)=\sigma\cdot n$，
$n=-e_1$ 时 $t=-\sigma_0e_1$；图注把两个端面写成同一个 $+\sigma_0e_1$，与正文的 $\sigma\cdot n$ 规则不一致。
从文字层看，图内两个标签也都是 $\sigma_0e_1$。

\newlbl 新稿图 8（p32）图注写：“右端面的外法向是 $+e_1$，面力 $t(+e_1)=+\sigma_0e_1$；左端面的外法向是
$-e_1$，面力 $t(-e_1)=-\sigma_0e_1$。两个向量都指向试件外侧，所以是‘拉’而不是‘压’。”图内标
$+\sigma_0e_1$、$-\sigma_0e_1$，灰色 $n$ 参考轴标 $30^\circ$。这一处同样是问题单 V1 之后改的，R1 的图 8 也
有端面符号与 $n$ 轴角度的不一致。

\bndlbl 本册只核对了图注与图内标签的文字层；箭头几何、箭头长度、遮挡与图注位置不由 DSH 目视。
round-04 的外部视觉审核已完成（见 \texttt{visual-review-round04.json}），本册不重复声称看图；若外部审核
对图 8 有具体意见，应查该报告而不是本册文字层。旧稿 $45^\circ$ 截面的最大剪应力 $\sigma_0/2$ 计算本身
正确；新稿把“最大剪应力在哪个角度”留到练习 P10，是编排差异，不是旧稿算错。

\clearpage
\hdr{C}{选读与主线：哪条路真的能走通}

本节先看最短路线（C1），再用旧稿相同主题的 §8.3、§9.4 对照新稿 §5.3、§6.4（C2、C3）。后两组不是为凑
组数：它们对应 round-04 接受范围内已保留的两处阅读桥梁，也是本次“只会简单线性代数”的读者最可能停下的
地方；旧稿在对应页有公式/结果，但缺这两步操作。原稿把旧稿 §4.6 的跳过提示说成“看不出可跳过”、把新稿
§1.4 的省略说成“只会在特征值一词停下”，这两条已撤回：旧稿 §4.6 段内已有“本讲义正文不用……那时再回来
补这半页即可”，新稿 §1.3 也已局部给出 $Tn=\lambda n$。

\cmpitem{C1}{旧稿的最短路线漏掉了 §7.1 要用的 §5.4}
\cropimg{old-p01-route}
\cropimg{old-p17-center}
\cropimg{new-p02-route}

\oldlbl 旧稿 p1 的“读法”给出最短路径 §1.2 → §2.1 → §3.3 → §4.3 → §5.1 → §7.1，并声明“这条路上不跳过
任何‘后面才用’的定义”。但 §7.1（p17）把 $r^2\hat1\cdot\omega$ 与 $(r\otimes r)\cdot\omega$ 都注为
“单位并矢的定义，§5.4”“投影并矢的作用，§5.4”；§5.4（p14）恰不在最短路径上。类似地，§2.1 的迹不变证明用
§1.3 的 $R$，而 §1.3 也被跳过。按字面走，读者在 p17 会撞到两个未读定义。

\newlbl 新稿 p2 的“最短主线”把 §0.3、§0.4、§2.2、§2.3、§2.5、§2.9 都列入，并说“跳过所有标（选读）的
段落即可；这些段落里没有后面要用的定义”。§2.5 正是单位并矢与投影并矢。R4 的阅读路径自查也逐条走了这条线，
记录了四处选读不喂后文。

\bndlbl 新稿主线仍省略 §1.4 的正式定义与 $\det(T-\lambda\hat1)=0$ 的求解构造；新稿 §1.3 已局部给出
$Tn=\lambda n$，§1.5 的谱定理陈述可独立使用。省略是事实，是否构成阅读阻断未试读，本册不再写“只会在
特征值一词回看一页”。另 §6.1 引用 §2.7 商法则，主线不列 §2.7。旧稿 §5.4 的内容本身正确，还配了图 2；
问题在“最短路线”的声明，不在这一节的内容。

\cmpitem{C2}{对称陀螺：旧稿用“标准振子方程”一句带过，新稿给局部验证与初值}
\cropimg{old-p20-symtop}
\cropimg{old-p21-symtop}
\cropimg{new-p29-symtop}
\cropimg{new-p29-bridge}

\oldlbl 旧稿 §8.3（p20--21）从 $\dot\omega_1=-\Omega\omega_2$、$\dot\omega_2=\Omega\omega_1$ 直接写
“这是标准的‘振子方程’，解为 $\omega_1=\omega_\perp\cos(\Omega t+\varphi_0)$、
$\omega_2=\omega_\perp\sin(\Omega t+\varphi_0)$”，随后只解释“绕体轴匀速转动”与几何图像。它没有把解
代回方程核对，也没有说明两个常数怎样由给定的初始横向矢量 $(\omega_1(0),\omega_2(0))$ 定出；读者若没学过
简谐振子，只能接受“标准”这个名称。

\newlbl 新稿 §5.3（p29）先不借助微分方程：把横向部分看成二维矢量
$w_\perp=(\omega_1,\omega_2)$，用 $\frac{d}{dt}(\omega_1,\omega_2)=\Omega(-\omega_2,\omega_1)$ 说明速度
垂直、大小 $|\Omega|\,|w_\perp|$；与 $(\omega_1,\omega_2)$ 点乘得 $\omega_1^2+\omega_2^2$ 为常数，再由
$\Omega$ 的正负说明旋转方向（$\Omega<0$ 时反向）。
然后写出 $\cos/\sin$ 解，并“直接代回可以当场核对”两条方程；给定 $w_\perp(0)$ 时给出
$\omega_\perp=\sqrt{\omega_1(0)^2+\omega_2(0)^2}$、$\cos\varphi_0=\omega_1(0)/\omega_\perp$、
$\sin\varphi_0=\omega_2(0)/\omega_\perp$，并举 $t=0$ 沿 $e_1$/$e_2$ 的例子。新稿还明说不需要预先学过
“简谐振子”。p30 另给 $A=B=4$、$C=5$、$\omega_3=6$ 的数值示范 $\Omega=1.5\,$rad/s、$T\approx4.19\,$s。

\bndlbl 新稿补的是局部验证与初值桥梁，不证明所有读者都能一次读懂；$\pi$ 相位与三角函数的读法仍按简单
线性代数之外的常识。旧稿的解式与 $\Omega$ 定义本身正确，缺的是这两步操作；本册不把“旧稿用‘标准’一词”
说成结论错误。学习者未试读。

\cmpitem{C3}{平衡方程：旧稿给公式，新稿把应力变化与净力接起来}
\cropimg{old-p22-balance}
\cropimg{new-p32-derive}
\cropimg{new-p32-example}

\oldlbl 旧稿 §9.4（p22）给出高斯定理的并矢版
$\int_S\sigma\cdot n\,dS=\int_V\nabla\cdot\sigma\,dV$、
$(\nabla\cdot\sigma)_k=\sum_i\partial_i\sigma_{ik}$，以及与体力平衡的两式
$\nabla\cdot\sigma+f=0$（静力学）、$\nabla\cdot\sigma+f=\rho a$（动力学），并写“写得开就是
$\partial_j\sigma_{ji}+f_i=0$——三个方程，正好对应三个方向的力平衡”。
它到高斯公式、平衡式与“写得开”一句为止，没有解释 $\partial_j$ 在做什么，也没有从一个随位置变化的
应力场算出净面力；均匀应力时为什么该项为零，也没有用例子说明。

\newlbl 新稿 §6.4（p32）先写“偏导在做什么”，解释 $\partial_j=\partial/\partial x_j$ 是“只让第 $j$ 个
坐标变，另外两个固定”，再取边长 $h$ 的小立方体：$+e_j$ 面与 $-e_j$ 面的面力差乘面积 $h^2$ 得净力
$\approx h^3\partial_j\sigma_{ij}$，除以 $h^3$ 就是单位体积净面力 $\partial_j\sigma_{ij}$；对 $j$ 求和即
$\nabla\cdot\sigma$。接着给 $\sigma(x_1)=\sigma_0x_1/L\,e_1\otimes e_1$ 的例子，算出
$(\nabla\cdot\sigma)_1=\sigma_0/L$，并说明 $+x_1$ 侧拉得更用力、静力学需要 $f=(-\sigma_0/L,0,0)$；再对照
“应力处处相同则 $\partial_1\sigma_{11}=0$、内部无净面力”的均匀单轴拉伸。

\bndlbl 新稿补的是操作意义与一个方向的展开，不是张量分析全章；严格高斯定理证明仍标选读。旧稿的公式、
静/动力学平衡式与“三个方程”说法正确；本册只指出旧稿没有把“应力变化$\rightarrow$净力”这一步写出来，
不把公式本身判错。学习者未试读。

\clearpage
\hdr{D}{未解释术语：力学词是否被当作已知}

\cmpitem{D1}{定点、质心与 $v=\omega\times r$（瞬时轴是新稿自己引入并解释的术语）}
\cropimg{old-p17-center}
\cropimg{new-p23-center}

\oldlbl 旧稿 §7.1（p17）开头：“对刚体，绕定点（或取质心为基点）转动时，任一点的速度是
$v=\omega\times r$。”定点、质心两个词没有解释，$v=\omega\times r$ 也没有理由；对只学过线性代数的读者，
这一段是术语加断言。旧稿这一段没有使用“瞬时轴”一词。

\newlbl 新稿 §4.1（p23）先解释 $v=\omega\times r$ 为什么成立（到瞬时轴距离为 $\rho$ 时速度大小
$\omega\rho$，方向垂直轴与 $r$ 张成的平面），再用旁注给“质心是质量加权平均位置
$R=\int r\,dm/m$”和“瞬时轴是某一瞬间速度为零的点的连线，角速度 $\omega$ 就沿它”。瞬时轴是新稿为解释
$v=\omega\times r$ 而引入的术语，新稿自己给了定义，不算旧稿漏讲。

\bndlbl 旁注只有一句一个定义，真正建立刚体运动学直觉仍要靠力学课；新稿也明说“这是刚体运动学的基本结论，
本稿从它出发”，没有假装证明。旧稿用定点/质心两个词而未解释，这一条仍成立；本册不把新稿自己引入的
“瞬时轴”倒算成旧稿的缺口。

\cmpitem{D2}{连续体与面力}
\cropimg{old-p21-continuum}
\cropimg{new-p30-continuum}

\oldlbl 旧稿 §9.1（p21）：“在连续体里过一点取一个截面……叫面力矢量。”没有解释连续体是什么，也没有说
面力是“哪一侧材料对哪一侧的作用”；随后直接引用柯西应力定理。

\newlbl 新稿 §6.1（p30）先写“把物体看成连续分布的物质，而不是一个个离散质点”，再给面力矢量的读法
“法向为 $n$ 的那一侧材料，对另一侧单位面积的作用力”，然后用四面体平衡给出线性性。

\bndlbl 连续体是力学模型，不是数学定义；新稿的一句解释也只够启动阅读。四面体推导是新增的中间步骤，本册
未逐行复核其物理假设（准静态、体积力作为高阶小量忽略）；这些属于原资料与力学课范围。

\clearpage
\hdr{E}{旧稿未被否定的部分：保留它已经成立的解释}

比较不等于把旧稿说成一无是处。下列内容在本册实际核对的两稿页面上都成立，且旧稿写得并不差：

\begin{itemize}
\item \textbf{§1.3 换基三步推导 + 数值核对}（旧 p3--4）。旧稿 §12 自己记录了“旧版用 $45^\circ$ 数值例子直接
跳到 $T'=R^{\mathsf T}TR$，2026-11-04 已改成三步推导 + 核对”；本册比较的 20261105 版已含这一修正。新稿
§0.4 保留同样的三步结构，并补了 $R^{\mathsf T}x$ 与“另一本书把 $R$ 的行当新基”的约定提醒。
\item \textbf{§2.1 迹不变的循环性证明}（旧 p5）：一行 $\mathrm{tr}(AB)=\mathrm{tr}(BA)$ 加 $RR^{\mathsf T}=\hat1$。
新稿 §0.5 保留同一证明。
\item \textbf{图 1 椭球与主轴}（旧 p8）：旧稿已画出 $\lambda=3/1$ 对应的半轴 $1/\sqrt3$ 与 1，并写“半轴
$=1/\sqrt\lambda$”。新稿图 1 保留这一读法，并加“与代入的 $u,v$ 逐项对应”。
\item \textbf{§7.3 对称轴判据、§7.4 平行轴定理的并矢版}（旧 p18）：新稿 §4.3/§4.4 保留同一结构和三角形
不等式验算。
\item \textbf{§4.6 上下标预告的跳过指引}（旧 p12）。旧稿在段内已明确写“这里给半页说明，本讲义正文不用”，
并说“那时再回来补这半页即可”；它已经提供了跳过与回补指引，只是没有把该小节标成“选读”。新稿 p20 的
差别是把“可跳过”写进标题、在段末给固定续读落点；这是导航标记的差别，不是旧稿没有提示。
\item \textbf{§12 缺口清单}（旧 p27--28）：旧稿逐条列出课件跳过的机制/压缩/拼接/图缺，并自记“以例代证”
的历史。这是有用的来源档案；新稿改用附录 C 与 \texttt{source-map.json}，形式不同，不能据此说旧稿没有
来源意识。
\end{itemize}

\srcsmall{这些不是“新稿更好”的证据。它们只说明旧稿并非全面失败，比较应针对具体页面的具体差别；旧稿保留的
优点不应为了让新稿显得更好而被删掉或贬低。}

\needspace{12\baselineskip}
\hdr{F}{仍未验证与不应声称的}

\begin{enumerate}
\item \textbf{学习者试读未发生。}不能声称已学会、已读懂或“效果更好”；本册只做页面与文字层比较。
\item \textbf{DSH 没有图像输入，本册仍需外部复核。}本册 28 处裁剪与本册最终页渲染已保存，需支持看图的
审核者检查字体、箭头、遮挡、图注位置；文字层检查不叫目视。新稿 round-04 的外部内容、数学、视觉审核
已经完成并记录在 \texttt{final-selection.json} 与 \texttt{visual-review-round04.json}，所以“DSH 没看图”
不等于“新稿没有外部视觉审核”；需要外部复核的是本修订版对比册本身。
\item \textbf{覆盖面有限。}本册只比较选出的 28 处裁剪；没有逐页审计旧稿 29 页或新稿 48 页，不能据此断言
“旧稿只有这些错”或“新稿没有其它错”。
\item \textbf{接受的边界。}新稿的 accepted 是“documented content, math and visual scope”内的接受；
\texttt{final-selection.json} 记录的外部 content/math/visual 审核已完成，\texttt{review\_limits} 明说未测
学习改善。接受不等于教学效果已验证。
\item \textbf{规则归因有限。}R2--R4 同时有具体问题单，无法把某处改动单独归因给 2.2.1/2.2.2 的某一条；本册
只把规则变化当作背景，不当作效果证据。
\item \textbf{交付。}本册 PDF 已按单份小于 25 MiB 检查；实际微信外发未在本会话执行（任务边界要求不外发），
交付动作由外部执行。
\item \textbf{旧稿版本。}旧稿封面印 20261103、文件版本为 20261105；本册按实际文件版本与物理页码区分，未另找
同名更早稿，也不据封面改称版本。
\end{enumerate}

\clearpage
\hdr{G}{本册的检查范围：文字层、页码、裁剪与编译}

以下检查由脚本执行，结果逐条写入 \texttt{verification.txt}；它们证明“本册引用的页面与文字确实来自这两份
PDF”“裁剪没有越界、锚字符串都在”“PDF 文字层与页码自洽”，不证明内容对学习者可读。

\begin{itemize}
\item 两份源 PDF 的页数、创建日期、SHA-256；旧稿 29 页、新稿 48 页与 \texttt{final-selection.json} 一致。
\item 28 处裁剪：矩形在页面内；裁剪文字层含每条 \texttt{require} 字符串；PNG 像素尺寸与 DPI 记录在案。
\item 本册 PDF：页数、每页文字层抽样、章节标题页位、裁剪图引用齐全、无未解析引用、日志无 Overfull/
Missing character。
\item 页边距：每页文字块在页面框内；图片不超出版心。
\item 文件大小与 SHA-256；最终每页渲染数量与 PDF 页数一致。
\end{itemize}

\warnbox{\textbf{未做。未看图。}上述检查全部基于 PDF 文字层、坐标与编译日志。字体是否缺字、箭头是否遮挡、
图注是否压线，仍需外部目视；本册不把文字层通过写成目视通过。新稿 round-04 的外部视觉审核已完成，
见 \texttt{final-selection.json} 与 \texttt{visual-review-round04.json}；这里说的是本册自身与裁剪未由
DSH 目视，本修订版对比册仍需外部复核。}

\clearpage
\hdr{附录}{裁剪清单、文件清单与复现}

\begingroup\small
\begin{longtable}{@{}p{3.55cm}p{1.05cm}p{0.95cm}p{3.9cm}p{3.15cm}@{}}
\thead{裁剪键} & \thead{来源} & \thead{页} & \thead{裁剪（pt）} & \thead{像素（200 dpi）}\\
\hline
\endhead
\texttt{old-p01-audience} & 旧稿 & 1 & 48,176--548,262 & 1390$\times$240 \\
\texttt{old-p01-route} & 旧稿 & 1 & 48,706--548,772 & 1390$\times$184 \\
\texttt{old-p11-summation} & 旧稿 & 11 & 48,300--548,458 & 1390$\times$440 \\
\texttt{old-p11-law} & 旧稿 & 11 & 48,676--548,792 & 1390$\times$323 \\
\texttt{old-p12-quotient} & 旧稿 & 12 & 48,326--548,470 & 1390$\times$401 \\
\texttt{old-p15-epsilon} & 旧稿 & 15 & 48,160--548,330 & 1390$\times$473 \\
\texttt{old-p15-note} & 旧稿 & 15 & 48,377--548,432 & 1390$\times$153 \\
\texttt{old-p17-center} & 旧稿 & 17 & 48,138--548,320 & 1390$\times$506 \\
\texttt{old-p20-symtop} & 旧稿 & 20 & 48,640--548,792 & 1390$\times$423 \\
\texttt{old-p21-symtop} & 旧稿 & 21 & 48,64--548,172 & 1390$\times$301 \\
\texttt{old-p21-continuum} & 旧稿 & 21 & 48,394--548,530 & 1390$\times$379 \\
\texttt{old-p22-fig5} & 旧稿 & 22 & 48,148--548,310 & 1390$\times$451 \\
\texttt{old-p22-balance} & 旧稿 & 22 & 48,320--548,500 & 1390$\times$501 \\
\texttt{old-p23-eig} & 旧稿 & 23 & 48,376--548,430 & 1390$\times$151 \\
\texttt{new-p01-assume} & 新稿 & 1 & 48,180--548,300 & 1390$\times$334 \\
\texttt{new-p02-actions} & 新稿 & 2 & 48,230--548,310 & 1390$\times$224 \\
\texttt{new-p02-route} & 新稿 & 2 & 48,146--548,216 & 1390$\times$195 \\
\texttt{new-p03-summation} & 新稿 & 3 & 48,104--548,192 & 1390$\times$246 \\
\texttt{new-p16-epsilon} & 新稿 & 16 & 48,684--548,794 & 1390$\times$306 \\
\texttt{new-p17-quotient} & 新稿 & 17 & 48,628--548,788 & 1390$\times$445 \\
\texttt{new-p23-center} & 新稿 & 23 & 48,410--548,556 & 1390$\times$407 \\
\texttt{new-p29-symtop} & 新稿 & 29 & 48,206--548,392 & 1390$\times$517 \\
\texttt{new-p29-bridge} & 新稿 & 29 & 48,395--548,685 & 1390$\times$806 \\
\texttt{new-p30-continuum} & 新稿 & 30 & 48,578--548,634 & 1390$\times$157 \\
\texttt{new-p32-derive} & 新稿 & 32 & 48,286--548,470 & 1390$\times$512 \\
\texttt{new-p32-example} & 新稿 & 32 & 48,478--548,676 & 1390$\times$551 \\
\texttt{new-p32-fig8} & 新稿 & 32 & 48,66--548,272 & 1390$\times$573 \\
\texttt{new-p38-p1} & 新稿 & 38 & 48,505--548,700 & 1390$\times$543 \\

\end{longtable}
\endgroup

{\raggedright
\textbf{文件清单。}\texttt{comparison.tex}（本册源文件）、\texttt{cropmeta.tex}/\texttt{croplist.tex}
（由脚本生成）、\texttt{crops.json}（裁剪定义）、\texttt{render\_crops.py}、\texttt{render\_pages.py}、
\texttt{verify\_old\_new\_math.py}、\texttt{check\_compare.py}、\texttt{check\_claims.py}、
\texttt{verification\_report.py}、\texttt{build.ps1}、\texttt{figs/}（28 张裁剪）、
\texttt{render/}（本册最终页渲染）、\texttt{verification.txt}、\texttt{issue-response.json}、
\texttt{delivery-manifest.json}、\texttt{render-crops.json}、\texttt{checks/}（编译与检查日志）。
两份源 PDF 只读引用 \texttt{inputs/old-handout.pdf} 与 \texttt{round-04/handout.pdf}，未复制进本目录。\par}

{\raggedright\sloppy
\textbf{复现。}\texttt{python -B render\_crops.py} → 生成 \texttt{figs/}、\texttt{cropmeta.tex}、
\texttt{croplist.tex}；两遍 \texttt{xelatex}；\texttt{python -B verify\_old\_new\_math.py}；
\texttt{python -B check\_compare.py --strict}；\texttt{python -B verification\_report.py}。也可直接运行
\texttt{build.ps1}。\par}

\end{document}
